فرض کنیم $C(\mathbb{R}, \mathbb{R})$ مجموعه تمام توابع پیوسته از $\mathbb{R}$ به $\mathbb{R}$ باشد. می‌خواهیم ثابت کنیم:
\[
|C(\mathbb{R}, \mathbb{R})| = |\mathbb{R}| = c
\]
که در آن $c = 2^{\aleph_0}$ نشان‌دهنده کاردینال مجموعه اعداد حقیقی است.

% \subsection*{۱. کران پایین: $|C(\mathbb{R}, \mathbb{R})| \ge |\mathbb{R}|$}
مجموعه تمام توابع ثابت را در نظر بگیرید. به ازای هر عدد حقیقی $k \in \mathbb{R}$، می‌توانیم تابع زیر را تعریف کنیم:
\[
f_k: \mathbb{R} \to \mathbb{R}, \quad f_k(x) = k
\]
هر تابع ثابت، در تمام دامنه خود پیوسته است. از آنجا که هر $k \in \mathbb{R}$ یک تابع پیوسته منحصربه‌فرد تولید می‌کند، این نگاشت یک‌به‌یک است. در نتیجه، کاردینال مجموعه توابع پیوسته حداقل برابر با کاردینال مجموعه اعداد حقیقی است:
\[
|C(\mathbb{R}, \mathbb{R})| \ge |\mathbb{R}| = c
\]

% \subsection*{۲. کران بالا: $|C(\mathbb{R}, \mathbb{R})| \le |\mathbb{R}|$}
یک ویژگی بنیادین توابع پیوسته این است که هر تابع پیوسته $f: \mathbb{R} \to \mathbb{R}$ به‌طور کامل با مقادیرش روی مجموعه اعداد گویا ($\mathbb{Q}$) تعیین می‌شود.

چون $\mathbb{Q}$ در $\mathbb{R}$ چگال است، برای هر عدد حقیقی $x \in \mathbb{R}$، دنباله‌ای از اعداد گویا مانند $\{q_n\}_{n=1}^{\infty}$ وجود دارد که به $x$ همگراست ($\lim_{n \to \infty} q_n = x$). بنا بر تعریف پیوستگی داریم:
\[
f(x) = f\left(\lim_{n \to \infty} q_n\right) = \lim_{n \to \infty} f(q_n)
\]
بنابراین، اگر دو تابع پیوسته $f$ و $g$ روی تمام نقاط گویا با هم برابر باشند (یعنی به ازای هر $q \in \mathbb{Q}$ داشته باشیم $f(q) = g(q)$)، باید روی تمام اعداد حقیقی نیز کاملاً با هم یکسان باشند.

این موضوع نشان می‌دهد که نگاشت محدودسازی $f \mapsto f|_\mathbb{Q}$ از مجموعه $C(\mathbb{R}, \mathbb{R})$ به مجموعه تمام توابع از $\mathbb{Q}$ به $\mathbb{R}$ (که با $\mathbb{R}^\mathbb{Q}$ نمایش داده می‌شود)، یک نگاشت یک‌به‌یک است. در نتیجه:
\[
|C(\mathbb{R}, \mathbb{R})| \le |\mathbb{R}^\mathbb{Q}|
\]

با استفاده از محاسبات کاردینالی و با توجه به اینکه $|\mathbb{Q}| = \aleph_0$ و $|\mathbb{R}| = 2^{\aleph_0}$ است، داریم:
\[
|\mathbb{R}^\mathbb{Q}| = |\mathbb{R}|^{|\mathbb{Q}|} = (2^{\aleph_0})^{\aleph_0} = 2^{\aleph_0 \cdot \aleph_0}
\]
از آنجا که حاصل‌ضرب دو کاردینال شمارای نامتناهی همچنان شماراست ($\aleph_0 \cdot \aleph_0 = \aleph_0$)، رابطه زیر برقرار خواهد بود:
\[
2^{\aleph_0 \cdot \aleph_0} = 2^{\aleph_0} = c = |\mathbb{R}|
\]
بنابراین:
\[
|C(\mathbb{R}, \mathbb{R})| \le |\mathbb{R}|
\]

با توجه به برقرار بودن هر دو نابرابری:
\begin{itemize}
    \item $|C(\mathbb{R}, \mathbb{R})| \ge |\mathbb{R}|$
    \item $|C(\mathbb{R}, \mathbb{R})| \le |\mathbb{R}|$
\end{itemize}
پس حکم زیر برقرار است:
\[
|C(\mathbb{R}, \mathbb{R})| = |\mathbb{R}| = c
\]
